% section: 解を組み合わせる見方

\section{解を組み合わせる}\label{節：解を組み合わせる}
更新式を満たす一つの\bookterm{sequence}{数列}を「解」と呼ぶ.
今度は\bookterm{initial}{初期条件}を固定せず,解全体を組み合わせから見よう.
この見方は隣接多項間・連立型・虚数解型,付録の\bookterm{matrix}{行列}に共通する.

\subsection{定数倍と足し算}\label{節：定数倍と足し算}
二つ前の項と同じになる
\[
  a_{n+2}=a_n\qquad(n\ge1)
\]
で,初めの二項を$(1,0)$,$(0,1)$とした解を$u_n,v_n$と書こう.
\[
  \begin{array}{c|rrrr}
    n&1&2&3&4\\ \hline
    u_n&1&0&1&0\\
    v_n&0&1&0&1
  \end{array}
\]
同じ番号の項をそれぞれ定数倍して足すと,$3u_n-2v_n$は$3,-2,3,-2,\ldots$となる.
これも同じ\bookterm{recurrence}{漸化式}を満たす.\footnote{項をすべて足して一つの数を作る総和とは異なる.番号$n$ごとに新しい項を作っている.}
定数$C,D$についても,
\[
  Cu_{n+2}+Dv_{n+2}=Cu_n+Dv_n
\]
だから,$Cu_n+Dv_n$は解である.
このように\bookterm{sequence}{数列}に定数を掛けて足したものを,それらの\textbf{\bookterm{combination}{線形結合}}という.

初めの二項は$(C,D)$となる.\bookterm{initial}{初期条件}を固定する前の解全体を考えている点に注意しよう.

\subsection{初期条件と解の数}\label{節：初期条件と解の数}
係数の重複,すべての解の表示,必要な数列の個数を順に確かめる.

\bookstep{係数の重複を確かめる}
今の$u_n,v_n$では,$Cu_n+Dv_n=0$がすべての$n$で成り立つなら,$n=1,2$から$C=D=0$となる.
初めの二項が$(C,D)$なので,係数が違えば数列も違う.
このような重複のない関係を,\textbf{\bookterm{independent}{一次独立}}という.\footnote{右辺の$0$は,すべての項が$0$の数列である.一つの番号で$0$になるだけでは足りない.例えば$u_2=0$だが,$u_n$は零数列ではない.}

一方,$a_{n+1}=2a_n$の解$u_n=2^{n-1}$と$v_n=2\cdot2^{n-1}$では,
\[
  Cu_n+Dv_n=(C+2D)2^{n-1}
\]
となり,係数は一意ではない.二つの非零の数列では,互いの定数倍かどうかで違いが分かる.

\bookstep{すべての解を組み立てる}
$a_{n+2}=a_n$で初めの二項が$(a_1,a_2)$なら,最初の$u_n,v_n$を使って
\[
  a_n=a_1u_n+a_2v_n
\]
と表せる.右辺は更新式と初めの二項を満たし,次の項が一意に決まるので,その後も一致する.
このように,\bookterm{independent}{一次独立}で,すべての解を\bookterm{combination}{線形結合}として表せる組を\textbf{\bookterm{basis}{基底}}という.
\bookterm{basis}{基底}を選べば,各解を表す係数も一意に決まる.

\bookstep{必要な数列の個数を数える}
今の例では二つで足り,一つでは初めの二項を自由に選べない.\bookterm{basis}{基底}の個数を\textbf{\bookterm{dimension}{次元}}という.
$a_{n+2}=a_n$は2\bookterm{dimension}{次元},$2^{n-1}$の定数倍で表せる$a_{n+1}=2a_n$は1\bookterm{dimension}{次元}である.

\bookstep{ほかの更新へ広げる}
一般に,$u_n,v_n$がともに
\[
  w_{n+2}+pw_{n+1}+qw_n=0
\]
の解なら,$b_n=Cu_n+Dv_n$について
\BookMathLayout{\[
    \begin{aligned}
      b_{n+2}+pb_{n+1}+qb_n
      &=C(u_{n+2}+pu_{n+1}+qu_n)\\
      &\quad+D(v_{n+2}+pv_{n+1}+qv_n)\\
      &=0
    \end{aligned}
  \]}{\[
    \begin{aligned}
      b_{n+2}+pb_{n+1}+qb_n\\
      {}=C(u_{n+2}+pu_{n+1}+qu_n)\\
      \quad+D(v_{n+2}+pv_{n+1}+qv_n)\\
      {}=0
    \end{aligned}
  \]}
となり,掛けて足す操作と更新が両立する.
このように,前の項に係数を掛けて足し,付加項を含まずに更新する式を\bookterm{homogeneous}{同次}\bookterm{linear}{線形}\bookterm{recurrence}{漸化式}という.\footnote{係数が$n$によって変わる場合も含む.用語の対応は\sectionref*{節：アフィン・同次・非同次}で確認できる.}
解を定数倍して足しても解になる性質を,\textbf{\bookterm{superposition}{重ね合わせの原理}}という.

\bookstep{数列の個数を増やす}
初期条件$(1,0)$と$(0,1)$から作った最初の解$u_n,v_n$に,$w_n=u_n+v_n$を加えよう.
どの二つも互いの定数倍ではないが,
\[
  u_n+v_n-w_n=0
\]
である.非零の係数$(1,1,-1)$でも零列を作れる.三つ以上では二つずつ比べるだけでは足りない.
数列の組を
\[
  u^{(1)}_n,\ldots,u^{(m)}_n
\]
と書く.上付きの$(1),\ldots,(m)$は,数列を区別する番号である.
\[
  C_1u^{(1)}_n+\cdots+C_mu^{(m)}_n=0
\]
がすべての$n$で成り立つのは$C_1=\cdots=C_m=0$の場合だけ,というのが一般の\bookterm{independent}{一次独立}の定義である.
\mbox{\bookterm{independent}{一次独立}}な解の組ですべての解を表せれば\bookterm{basis}{基底}となり,その個数が\bookterm{dimension}{次元}となる.

この解全体は零\bookterm{sequence}{数列}を含み,足し算や定数倍でも解のままである.\footnote{こうした操作のできる集合を\textbf{\bookterm{vector-space}{ベクトル空間}}という.}
ここでは係数を実数としている.\footnote{後の虚数解型で複素数を使うときも,係数を複素数にして同じ計算と判定を使える.どちらの係数で考えるかは区別する.}

\bookterm{dimension}{次元}は,\bookterm{initial}{初期条件}として自由に選べる数値の個数とつながっている.\footnote{\bookterm{nonlinear}{非線形}の式や,途中で更新が定義されなくなる式に,この結論をそのまま当てはめないようにしよう.}
正の整数$k$に対し,すべての$n\ge1$で係数が定義された\bookterm{homogeneous}{同次}\bookterm{linear}{線形}\bookterm{recurrence}{漸化式}
\[
  a_{n+k}=\sum_{j=0}^{k-1}c_j(n)a_{n+j}
\]
では,$c_j(n)$は既知の係数,$j$は和の番号である.初めの$k$項から次が一意に決まる.
初めの$k$項のうち一つだけを$1$,残りを$0$とした$k$個の解を作れば,
それらの\bookterm{combination}{線形結合}でどんな\bookterm{initial}{初期条件}にも合わせられる.
係数の一意性も初めの$k$項を比べれば分かるので,解全体は$k$\bookterm{dimension}{次元}である.
\bookterm{initial}{初期条件}を指定した後は,その中から一つの解を選んでいる.

\subsection{解の差を取る}\label{節：解の差を取る}
$a_{n+1}=2a_n+3$の二つの解を足すと付加項は$6$になる.\bookterm{inhomogeneous}{非同次}では任意の\bookterm{combination}{線形結合}は解にならない.
一方,同じ更新をする二つの解$a_n,g_n$を引けば,
\[
  (a_{n+1}-g_{n+1})=2(a_n-g_n)
\]
となる.共通の付加項が消え,\bookterm{homogeneous}{同次}の式が残る.
この例では$g_n=-3$を選べるから,
\[
  a_n=-3+C2^{n-1}
\]
と全体を表せる.初項から$C=a_1+3$を決める.
\sectionref*{節：補助数列を作る}で引いた量は,既知の解と\bookterm{homogeneous}{同次}の解に分ける役割もあった.

この見方は,後の節で次のように使う.
\begin{itemize}
  \item \sectionref*{節：階差等比複合型}では,同じ付加項をもつ解を一つ見つけ,引き算で\bookterm{homogeneous}{同次}に直す.
  \item \sectionref*{節：隣接3項型}では,\bookterm{root}{特性根}から独立な二つの解を作り,\bookterm{initial}{初期条件}を合わせる.
    重解のときは同じ冪を二つ並べても足りないので,別の形の解を探す.
  \item \sectionref*{節：連立型}では,二つの量の和と差を使って更新を分ける.
    \hyperref[節：特性根が虚数の場合]{特性根が虚数の場合}では,複素数係数で使う冪の\bookterm{basis}{基底}を,
    実数係数で使える正弦・余弦の\bookterm{basis}{基底}へ直す.
\end{itemize}
「定数を掛けて足す」操作には,解全体を組み立てる役割と,更新を簡単にする\bookterm{auxiliary}{補助数列}を作る役割がある.
付録の\sectionref*{節：行列の基本}では,この二つの役割をベクトルと\bookterm{matrix}{行列}から見直す.
